\chapter{The Matrix-Free Operator}
\label{ch:matrixfree}

\chapterorient{Assembly (\cref{ch:assembly}) built a matrix. This chapter
un-builds it. We will see that a finite-element operator does not have to exist
as a stored array of numbers at all: the only thing a solver ever asks of an
operator is its \emph{action}, the map $\vv \mapsto \mat{A}\vv$, and that action
can be computed on the fly as a short chain of element-local tensor
contractions---gather, evaluate, weight, integrate, scatter---without ever
forming or storing $\mat{A}$. This is the \emph{matrix-free} operator. It is the
form that costs the least memory at high polynomial order, and---because every
stage is a contraction over named tensor axes---it is exactly the form a GPU or
other tensor accelerator wants to run. By the end you will be able to trace a
vector through all five stages, recover the same answer the assembled matrix
would have given, and see why the innermost kernel is the place the book's
``implementation view'' finally touches the silicon.}

\section{What a solver actually asks for}
\label{sec:mf-whatask}

Open any of the iterative solvers in \cref{part:framework}---conjugate
gradients (\cref{ch:cg}), GMRES (\cref{ch:gmres}), the Lanczos and Arnoldi
recurrences behind the eigensolver (\cref{ch:krylovschur})---and look at what
they do with the operator $\mat{A}$. They form inner products, they orthogonalize,
they take linear combinations of vectors. But the \emph{only} way any of them
touches $\mat{A}$ is through one operation: given a vector $\vv$, produce the
vector $\mat{A}\vv$. They never read an individual entry $A_{ij}$. They never ask
for a row, a column, or the sparsity pattern. A Krylov method is, by
construction, a procedure that explores the subspace
$\{\vv,\,\mat{A}\vv,\,\mat{A}^2\vv,\dots\}$, and to walk that subspace it needs
nothing but repeated application of the \emph{action}
\[
  \vv \;\longmapsto\; \mat{A}\vv .
\]

This is worth pausing on, because it inverts the picture most readers carry from
a first linear-algebra course. There, a matrix \emph{is} its array of entries;
``apply the matrix'' means ``do the row-times-column sums.'' But an
iterative solver does not need the entries---it needs a black box that, fed
$\vv$, returns $\mat{A}\vv$. How that box computes its answer is its own affair.
Assembly (\cref{ch:assembly}) is one way to build such a box: store every entry,
then multiply. This chapter is the other way.

\begin{keyidea}
A Krylov solver is a consumer of the operator's \emph{action}, $\vv\mapsto
\mat{A}\vv$, and of nothing else. It never inspects the entries of $\mat{A}$. So
any object that can answer ``given $\vv$, here is $\mat{A}\vv$'' is, to the
solver, a perfectly good operator---whether or not a matrix was ever built. This
single fact is the licence for everything in this chapter.
\end{keyidea}

\subsection{Why ``store the matrix'' gets expensive}
\label{sec:mf-cost}

If the action is all we need, why ever store the matrix? For low polynomial
order, storing it is cheap and fast, and assembly (\cref{ch:assembly}) is the
right tool. The trouble begins at \emph{high order}, which is exactly where
finite-element accuracy is won most cheaply.

Recall the element-local stiffness or mass matrix from \cref{ch:assembly}: on one
element with $L$ local degrees of freedom, the element matrix $\mat{A}^e$ is
$L\times L$, holding $L^2$ numbers. For a tensor-product element of polynomial
degree $p$ in $d$ spatial dimensions, the number of local degrees of freedom
grows as
\[
  L \;=\; (p+1)^d ,
\]
so the dense element matrix holds
\[
  L^2 \;=\; (p+1)^{2d}
\]
numbers. In three dimensions ($d=3$) that is $(p+1)^6$ entries \emph{per element}.
At $p=1$ this is a tidy $2^6 = 64$; at $p=4$ it is $5^6 = \num{15625}$; at $p=8$
it is $9^6 = \num{531441}$---half a million numbers for a single element. Storing
the assembled operator means storing something on this order, multiplied across
every element of the mesh. The memory cost of the \emph{stored} operator scales
as $O(p^{2d})$ per element, and on modern accelerators memory---not
arithmetic---is the scarce resource.

\Cref{fig:mf-cost} plots the contrast. The assembled operator's storage climbs
as $(p+1)^{2d}$; the matrix-free operator stores only the tabulated
one-dimensional basis and the per-quadrature-point geometry factor, a cost that
grows far more gently. Past a modest order---in 3-D, around $p=2$ to $4$
depending on the architecture---the matrix-free form is not merely lighter, it is
the only form that fits.

\begin{figure}[t]
\centering
\begin{tikzpicture}
\begin{axis}[
  width=0.86\linewidth, height=58mm,
  xlabel={polynomial order $p$},
  ylabel={numbers stored per element},
  ymode=log,
  xmin=1, xmax=8, ymin=1, ymax=1e6,
  xtick={1,2,3,4,5,6,7,8},
  ytick={1e0,1e2,1e4,1e6},
  legend pos=north west,
  legend cell align=left,
  grid=both, grid style={simGray!18},
  tick label style={font=\footnotesize},
  label style={font=\small},
  legend style={font=\footnotesize},
]
% assembled: (p+1)^6  (3-D, d=3, L^2)
\addplot[simRust, very thick, mark=*, mark size=1.6pt]
  coordinates {(1,64)(2,729)(3,4096)(4,15625)(5,46656)(6,117649)(7,262144)(8,531441)};
\addlegendentry{assembled $\;(p{+}1)^{6}$}
% matrix-free: tabulated 1-D basis + geom factor ~ O(p) +small const
% store ~ (p+1)^2 (1-D basis) + a few geom comps per quad pt; use (p+1)^2 + 6(p+1)
\addplot[simBlue, very thick, mark=square*, mark size=1.6pt]
  coordinates {(1,16)(2,27)(3,40)(4,55)(5,72)(6,91)(7,112)(8,135)};
\addlegendentry{matrix-free (1-D basis $+$ geom)}
\end{axis}
\end{tikzpicture}
\caption[Memory: assembled versus matrix-free]{Numbers stored per element, in
3-D, as polynomial order $p$ rises (log scale). The \emph{assembled} operator
holds the dense element matrix, $(p+1)^{2d}=(p+1)^6$ entries, and climbs
steeply. The \emph{matrix-free} operator stores only the tabulated
one-dimensional basis tables and a handful of geometry components per quadrature
point, growing roughly linearly. The two curves cross almost immediately;
beyond low order, matrix-free is the only form that fits in accelerator memory.}
\label{fig:mf-cost}
\end{figure}

\begin{physicsbox}
The arithmetic the operator performs is the same in both forms---a weighted sum
of basis-function interactions over each element, the discretization of the weak
form $a(u,v)=\inner{Q\diffop u}{\diffop v}$ from \cref{ch:galerkin}. What
differs is \emph{bookkeeping}: the assembled form pre-computes and stores all
those interactions as a matrix; the matrix-free form recomputes them, cheaply,
each time the action is needed. We are trading a little extra arithmetic
(recompute on demand) for a large saving in memory (store nothing). On hardware
where memory bandwidth is the bottleneck and arithmetic is nearly free---the
defining shape of a modern GPU---this is overwhelmingly the right trade.
\end{physicsbox}

\section{The action as a chain of contractions}
\label{sec:mf-chain}

So we want to compute $\mat{A}\vv$ without $\mat{A}$. The route is to go back to
where the matrix came from. In \cref{ch:assembly} we built $\mat{A}$ by
\emph{assembling} element contributions: each element computed a local operator
from (i) restricting the global vector to the element, (ii) evaluating the
trial field and its derivatives at the element's quadrature points, (iii)
weighting those values by the geometry and material at each point, and (iv)
integrating back to the element's degrees of freedom; then (v) the element
results were summed into the global vector. Assembly did this \emph{once}, for
every basis vector, and stored the resulting matrix. Matrix-free does exactly the
same five operations, but \emph{on the actual input vector $\vv$}, every time the
action is requested---and stores nothing.

Each of the five operations is a linear map, so the whole action is their
composition. Reading right to left in the order they act on $\vv$:
\begin{equation}
  \mat{A}\,\vv
   \;=\;
   \underbrace{G^{\mathsf T}}_{\text{scatter}}\;
   \underbrace{B_{\diffop}^{\mathsf T}}_{\text{integrate}}\;
   \underbrace{D}_{\text{weight}}\;
   \underbrace{B_{\diffop}}_{\text{evaluate}}\;
   \underbrace{G}_{\text{gather}}\;
   \vv .
  \label{eq:mf-chain}
\end{equation}
This factorization is the heart of the chapter. It says the finite-element
operator is a \emph{sandwich}: a gather $G$ and its transpose scatter
$G^{\mathsf T}$ on the outside; a basis evaluation $B_{\diffop}$ and its transpose
$B_{\diffop}^{\mathsf T}$ inside that; and at the very center a single pointwise
multiply $D$. The symmetry is not a coincidence---it is the discrete shadow of
the weak form's own symmetry, $a(u,v)=a(v,u)$ for the self-adjoint terms, with
$G$ and $B$ carrying the trial side and their transposes carrying the test side
(\cref{ch:galerkin}).

\begin{notationbox}
The subscript $\diffop$ on $B_{\diffop}$ records \emph{which} differential
operator the weak-form term uses---value (identity), gradient, curl, or
divergence---the four members of the de Rham family from \cref{ch:derham}. A
mass term ($\diffop = \mathrm{Id}$) evaluates basis \emph{values}; a stiffness
term ($\diffop=\Grad$) evaluates basis \emph{gradients}; a curl--curl term
evaluates \emph{curls}. The choice of $\diffop$ is the only thing that changes
between the operators of the six analyses (\cref{ch:targets}); the five-stage
skeleton is identical for all of them.
\end{notationbox}

\subsection{The five stages and their tensor shapes}
\label{sec:mf-stages}

Each stage carries a tensor from one named shape to the next. The axes are the
element-local family we will fix carefully in \cref{sec:mf-anatomy}; for now read
$N$ as the global degree-of-freedom count, $E$ as the number of elements, $L$ as
local degrees of freedom per element, $P$ as quadrature points per element, and
$C$ as the number of value (or derivative) components carried at each quadrature
point. \Cref{fig:mf-pipeline} is the hero diagram of the chapter; the prose
below walks its five boxes left to right.

\begin{figure}[p]
\centering
\resizebox{\linewidth}{!}{%
\begin{tikzpicture}[node distance=8mm and 9mm]
  % input
  \node[data, text width=18mm] (in) {global vector\\\texttt{Tensor[N]}};
  % G
  \node[stage, right=of in, text width=20mm] (G) {$G$\\\footnotesize gather\\(element restrict)};
  % B
  \node[stage, right=of G, text width=20mm] (B) {$B_{\diffop}$\\\footnotesize evaluate\\(basis apply)};
  % D  -- the extern kernel center
  \node[extern, right=of B, text width=20mm] (D) {$D$\\\footnotesize weight\\\texttt{\#extern}};
  % Bt
  \node[stage, right=of D, text width=20mm] (Bt) {$B_{\diffop}^{\mathsf T}$\\\footnotesize integrate};
  % Gt
  \node[stage, right=of Bt, text width=20mm] (Gt) {$G^{\mathsf T}$\\\footnotesize scatter-add};
  % output
  \node[data, right=of Gt, text width=18mm] (out) {global vector\\\texttt{Tensor[N]}};

  \draw[flow] (in) -- (G);
  \draw[flow] (G) -- (B);
  \draw[flow] (B) -- (D);
  \draw[flow] (D) -- (Bt);
  \draw[flow] (Bt) -- (Gt);
  \draw[flow] (Gt) -- (out);

  % shape annotations between stages (below the arrows)
  \node[font=\scriptsize\ttfamily, simGray, below=2.5mm of G.south] {[E, L]};
  \node[font=\scriptsize\ttfamily, simGray, below=2.5mm of B.south] {[E, P, C]};
  \node[font=\scriptsize\ttfamily, simGray, below=2.5mm of D.south] {[E, P, C]};
  \node[font=\scriptsize\ttfamily, simGray, below=2.5mm of Bt.south] {[E, L]};
  \node[font=\scriptsize\ttfamily, simBlue, below=3mm of in.south, align=center] {action\\input};
  \node[font=\scriptsize\ttfamily, simBlue, below=3mm of out.south, align=center] {action\\output};

  % geom_data feeding D from below
  \node[data, fill=simBgGold, draw=simGold, below=12mm of D, text width=26mm]
    (geom) {\footnotesize\itshape geom\_data \texttt{[E, P, G]}\\(build-time factor)};
  \draw[flow, simGold] (geom) -- (D);

  % brackets: outer = gather/scatter sandwich, inner = eval/integrate
  \draw[decorate, decoration={brace, amplitude=4pt}, simRust]
    ([yshift=7mm]G.north west) -- ([yshift=7mm]Gt.north east)
    node[midway, above=4pt, font=\scriptsize, simRust]
    {element-local: independent, parallel over $E$};
\end{tikzpicture}%
}
\caption[The five-stage matrix-free pipeline]{The matrix-free operator
$\mat{A}\vv = G^{\mathsf T}B_{\diffop}^{\mathsf T}\,D\,B_{\diffop}\,G\,\vv$ as a
dataflow pipeline, with the tensor shape annotated on each wire. A global vector
\lstinline[style=l4style]|Tensor[N]| enters; $G$ \emph{gathers} it to per-element
local tensors \lstinline[style=l4style]|[E, L]|; $B_{\diffop}$ \emph{evaluates}
the field (or its gradient/curl) at the \lstinline[style=l4style]|[E, P, C]|
quadrature points; the center stage $D$ \emph{weights} each point by the
precomputed \texttt{geom\_data} (gold, fed from below; the
\lstinline[style=l4style]|#extern| innermost kernel, \cref{sec:mf-impl-view});
$B_{\diffop}^{\mathsf T}$ \emph{integrates} back to local dofs; and $G^{\mathsf T}$
\emph{scatters} the element results into a fresh global vector. The whole middle
of the pipeline (rust brace) is independent across elements---the source of all
the parallelism.}
\label{fig:mf-pipeline}
\end{figure}

\paragraph{Stage 1 --- $G$, the gather (element restriction).}
The input is the global vector $\vv$, a \lstinline[style=l4style]|Tensor[N]| with
one entry per global degree of freedom. The gather $G$ reads, for each element,
the values of the degrees of freedom that element touches, copying them into a
per-element block. The output is a \lstinline[style=l4style]|Tensor[E, L]|: $E$
elements, each holding its $L$ local degrees of freedom. No arithmetic happens
here---$G$ is pure copying, a Boolean selection that fans the shared global
vector out into (generally overlapping) element-local pieces. A degree of freedom
sitting on a shared element boundary is copied into \emph{every} element that
touches it; that duplication is intentional and is undone, by summation, in the
final stage.

\paragraph{Stage 2 --- $B_{\diffop}$, the basis evaluation.}
Now each element holds its local coefficients. To apply the weak form we need the
\emph{field} those coefficients represent---and its derivative---sampled at the
element's quadrature points. That is exactly what the tabulated basis gives us
(\cref{ch:refelement}): a small dense matrix whose entries are the basis
functions (or their derivatives) evaluated at the quadrature nodes. $B_{\diffop}$
contracts each element's \lstinline[style=l4style]|[L]| coefficient vector against
that matrix, producing \lstinline[style=l4style]|[P, C]| values---$C$ components
($C=1$ for a scalar value, $C=d$ for a gradient) at each of $P$ quadrature points.
Stacked over elements, the output is a \lstinline[style=l4style]|Tensor[E, P, C]|.
The mode---values, gradient, curl, or divergence---is the $\diffop$ the term
asked for.

\paragraph{Stage 3 --- $D$, the pointwise quadrature weight.}
This is the center of the sandwich, and it is the simplest stage of all: at each
quadrature point, multiply the evaluated value by a single precomputed factor and
move on. The factor bundles three things---the material coefficient $Q$ (the
$\varepsilon$, $\mu^{-1}$, $\dots$ of the weak form), the geometry metric derived
from the element's Jacobian, and the quadrature weight $w$---into one number (or,
for vector terms, one small matrix) per point, stored in a tensor
\texttt{geom\_data} of shape \lstinline[style=l4style]|Tensor[E, P, G]|. There is
\emph{no coupling between points}: the output at point $(e,p)$ depends only on the
input at $(e,p)$ and on \texttt{geom\_data}$[e,p]$. The shape is unchanged,
\lstinline[style=l4style]|[E, P, C]| in and out (we write $C'$ for the output
component count, usually equal to $C$). \Cref{fig:mf-diagonal} shows why this is
the \emph{embarrassingly parallel} diagonal of the whole operator.

\paragraph{Stage 4 --- $B_{\diffop}^{\mathsf T}$, integrate back to dofs.}
The weighted quadrature-point values must now be turned back into element-local
degree-of-freedom contributions---this is the discrete integral $\int (\cdot)\,
\phi_i$ against each test basis function $\phi_i$. It is the exact transpose of
stage 2: the same tabulated basis matrix, applied on the other side. The shape
goes back from \lstinline[style=l4style]|[E, P, C]| to
\lstinline[style=l4style]|[E, L]|.

\paragraph{Stage 5 --- $G^{\mathsf T}$, the scatter-add (assembly).}
Finally, the per-element \lstinline[style=l4style]|[E, L]| contributions must be
returned to the single global vector. $G^{\mathsf T}$ is the transpose of the
gather: where $G$ copied a shared dof out to several elements, $G^{\mathsf T}$
\emph{sums} the several element contributions back into that one global slot.
This summation is precisely the finite-element ``assembly'' step
(\cref{ch:assembly})---here applied to a vector rather than a matrix. The output
is a fresh \lstinline[style=l4style]|Tensor[N]|: the action $\mat{A}\vv$, complete.

\begin{keyidea}
The matrix-free operator is a chain of tensor contractions over named axes---no
sparse-matrix data structure anywhere. \emph{Gather} ($N\!\to\!E,L$),
\emph{evaluate} ($E,L\!\to\!E,P,C$), \emph{weight} (pointwise on $E,P,C$),
\emph{integrate} ($E,P,C\!\to\!E,L$), \emph{scatter} ($E,L\!\to\!N$). The matrix
$\mat{A}$ that assembly would have stored is recovered, on demand, as the
composition $G^{\mathsf T}B_{\diffop}^{\mathsf T}DB_{\diffop}G$---and never
materialized.
\end{keyidea}

\subsection{The chain in pseudocode}
\label{sec:mf-pseudo}

In the calculus of \cref{part:framework}, the action reads as a single
composition of the five stage functions. Each stage is a pure tensor-to-tensor
map; threading a vector through them is exactly the value-threading picture of
\cref{ch:bigidea}.

\begin{lstlisting}[style=l4style]
-- the five element-local stages, as pure tensor maps
gather    :: ElemRestriction -> Tensor[N]        -> Tensor[E, L]      -- G
evaluate  :: BasisMode -> Basis -> Tensor[E, L]  -> Tensor[E, P, C]  -- B_D
weight    :: GeomData -> Tensor[E, P, C]         -> Tensor[E, P, C]  -- D
integrate :: BasisMode -> Basis -> Tensor[E, P, C] -> Tensor[E, L]   -- B_D^T
scatter   :: ElemRestriction -> Tensor[E, L]     -> Tensor[N]        -- G^T

-- the matrix-free action A v = G^T . B_D^T . D . B_D . G  v
apply_matrixfree :: Operator -> Tensor[N] -> Tensor[N]
apply_matrixfree op v =
  scatter   op.restr           $   -- G^T  : [E, L] -> [N]
  integrate op.mode op.basis   $   -- B_D^T: [E, P, C] -> [E, L]
  weight    op.geom            $   -- D    : pointwise on [E, P, C]
  evaluate  op.mode op.basis   $   -- B_D  : [E, L] -> [E, P, C]
  gather    op.restr v             -- G    : [N] -> [E, L]
\end{lstlisting}

\noindent Read it bottom to top, the order the data flows: \texttt{gather} the
global vector to element blocks, \texttt{evaluate} the basis, \texttt{weight}
pointwise, \texttt{integrate} back, \texttt{scatter} to a fresh global vector.
The operator value \texttt{op} carries the three \emph{build-time} pieces---the
restriction index map \texttt{op.restr}, the tabulated \texttt{op.basis}, and the
geometry factor \texttt{op.geom}---which we turn to next.

\section{The element-local tensor and its axes}
\label{sec:mf-anatomy}

The five stages move data between three rank-structured tensors. Naming their
axes precisely is what lets us reason about shapes the way \cref{ch:tensors}
taught---and it is what makes the whole pipeline legible to a tensor backend.
\Cref{fig:mf-anatomy} is the anatomy chart; \cref{tab:mf-axes} is the legend.

The crucial conceptual move is that matrix-free FE assembly works over a
\emph{richer} shape than the flat global vector. The solver sees a flat
\lstinline[style=l4style]|Tensor[N]|---a rank-1 list of $N$ degrees of freedom.
But the moment we gather to elements, we acquire an explicit \emph{element axis}
$E$ and the per-element structure beneath it. The flat global vector and the
element-local tensors are two different vocabularies, and the gather/scatter pair
$G$, $G^{\mathsf T}$ is exactly the boundary between them.

\begin{table}[t]
\centering
\caption[The element-local tensor axes]{The five named axes of the element-local
tensor family. The \emph{stratum} column records whether an axis's extent is
fixed once per $(\text{mesh},\text{order})$ in a setup pass (build-time) or
indexes a value produced afresh on every operator application (run-time)---the
distinction \cref{sec:mf-strata} makes load-bearing.}
\label{tab:mf-axes}
\small
\setlength{\tabcolsep}{6pt}
\renewcommand{\arraystretch}{1.25}
\begin{tabular}{@{}clll@{}}
\toprule
axis & name & meaning & stratum \\
\midrule
$E$ & element count & number of mesh elements assembled over & build-time \\
$L$ & local dofs & dofs of the basis local to one element, $(p{+}1)^d$ & build-time \\
$P$ & quad points & quadrature points per element & build-time \\
$C$ & components & value/derivative components at a quad point & run-time \\
$G$ & geom components & precomputed metric storage per quad point & build-time \\
\bottomrule
\end{tabular}
\end{table}

\begin{figure}[t]
\centering
\begin{tikzpicture}[scale=1.0]
  % three stacked tensors with axis labels
  % --- [E, L] ---
  \begin{scope}[shift={(-4.4,0)}]
    \foreach \k in {0,1,2}{
      \begin{scope}[shift={(0.12*\k,0.09*\k)}]
        \fill[simBgBlue, draw=simBlue] (-0.7,-0.5) rectangle (0.7,0.5);
        \foreach \y in {-0.25,0,0.25}{\draw[simBlue!40] (-0.7,\y)--(0.7,\y);}
      \end{scope}}
    \node[font=\scriptsize, simInk] at (-1.05,0.0) {$L$};
    \node[font=\scriptsize, simRust] at (0.55,0.75) {$E$};
    \node[font=\scriptsize\ttfamily, simInk] at (0.1,-1.0) {[E, L]};
    \node[font=\scriptsize\itshape, simGray, text width=24mm, align=center]
      at (0.1,-1.5) {per-element local dofs};
  \end{scope}
  \node[font=\large] at (-2.5,0) {$\xrightarrow{\,B_{\diffop}\,}$};
  % --- [E, P, C] ---
  \begin{scope}[shift={(-0.4,0)}]
    \foreach \k in {0,1,2}{
      \begin{scope}[shift={(0.12*\k,0.09*\k)}]
        \fill[simBgGreen, draw=simGreen] (-0.75,-0.5) rectangle (0.75,0.5);
        \foreach \y in {-0.17,0.17}{\draw[simGreen!40] (-0.75,\y)--(0.75,\y);}
        \foreach \x in {-0.25,0.25}{\draw[simGreen!40] (\x,-0.5)--(\x,0.5);}
      \end{scope}}
    \node[font=\scriptsize, simInk] at (-1.1,0.0) {$P$};
    \node[font=\scriptsize, simInk] at (0.0,-0.85) {$C$};
    \node[font=\scriptsize, simRust] at (0.6,0.75) {$E$};
    \node[font=\scriptsize\ttfamily, simInk] at (0.05,-1.15) {[E, P, C]};
    \node[font=\scriptsize\itshape, simGray, text width=26mm, align=center]
      at (0.05,-1.65) {values at quad points};
  \end{scope}
  \node[font=\large] at (1.6,0) {$\odot$};
  % --- [E, P, G]  geom_data ---
  \begin{scope}[shift={(3.7,0)}]
    \foreach \k in {0,1,2}{
      \begin{scope}[shift={(0.12*\k,0.09*\k)}]
        \fill[simBgGold, draw=simGold] (-0.75,-0.5) rectangle (0.75,0.5);
        \foreach \y in {-0.17,0.17}{\draw[simGold!50] (-0.75,\y)--(0.75,\y);}
        \foreach \x in {-0.25,0.25}{\draw[simGold!50] (\x,-0.5)--(\x,0.5);}
      \end{scope}}
    \node[font=\scriptsize, simInk] at (-1.1,0.0) {$P$};
    \node[font=\scriptsize, simInk] at (0.0,-0.85) {$G$};
    \node[font=\scriptsize, simRust] at (0.6,0.75) {$E$};
    \node[font=\scriptsize\ttfamily, simInk] at (0.05,-1.15) {[E, P, G]};
    \node[font=\scriptsize\itshape, simGray, text width=26mm, align=center]
      at (0.05,-1.65) {geom\_data (build-time)};
  \end{scope}
\end{tikzpicture}
\caption[Anatomy of the element-local tensor]{The element-local tensor family.
The gather produces \lstinline[style=l4style]|[E, L]| (blue): an element axis $E$
over per-element local-dof blocks of length $L$. Basis evaluation
$B_{\diffop}$ maps it to \lstinline[style=l4style]|[E, P, C]| (green): $C$
components at each of $P$ quadrature points, per element. The pointwise weight
$D$ multiplies it, point by point ($\odot$), against the precomputed geometry
carrier \lstinline[style=l4style]|[E, P, G]| (gold). The element axis $E$ (rust)
runs through all three---the operator is block-diagonal in $E$, which is to say
embarrassingly parallel over elements.}
\label{fig:mf-anatomy}
\end{figure}

\begin{notationbox}
A naming hazard worth flagging once. In the underlying libCEED library the basis
constructor's parameter named \texttt{P} counts \emph{nodes per element} (our
$L$) and its \texttt{Q} counts \emph{quadrature points} (our $P$). This book uses
$L$ for local dofs and $P$ for quadrature points consistently; if you read the
library source, expect the letters to swap. The mapping is exact---only the
labels differ.
\end{notationbox}

\subsection{Why every stage is block-diagonal in \texorpdfstring{$E$}{E}}
\label{sec:mf-blockdiag}

Look again at \cref{fig:mf-anatomy}: the element axis $E$ threads through all
three tensors, and through stages 2, 3, and 4 of the pipeline. None of those
stages couples one element to another. Basis evaluation applies the same small
matrix \emph{independently} within each element; the pointwise weight touches one
point at a time; integration is again per-element. In the language of
\cref{ch:tensors}, stages 2--4 are \emph{block-diagonal} in $E$: a batched
operation, the same computation broadcast across the $E$ slices.

All the inter-element coupling---the only place where one element's data meets
another's---lives in the gather $G$ and the scatter $G^{\mathsf T}$, at the two
ends of the pipeline. This is a clean and consequential separation. The
arithmetic-heavy middle is perfectly parallel; the communication-heavy ends are
pure data movement with no arithmetic. It is exactly the structure a tensor
accelerator is built to exploit: map a batched contraction over a large axis,
with a gather in front and a scatter behind.

\section{Build-time versus run-time: where the geometry lives}
\label{sec:mf-strata}

We have been quietly carrying a tensor we have not built: \texttt{geom\_data},
the \lstinline[style=l4style]|Tensor[E, P, G]| that the weight stage $D$
multiplies against. Where does it come from, and---more importantly---\emph{when}?
The answer is a recurring theme of this book, large enough that it has its own
chapter ahead (\cref{ch:strata}): the operator's data divides into two
\emph{strata} by lifetime.

\begin{description}[leftmargin=1.4em, style=nextline]
\item[Build-time (setup) stratum.] Data fixed once per
$(\text{mesh},\text{FE order},\text{quadrature rule})$ and reused across every
application of the operator. This is the restriction index map (which dof goes to
which element), the tabulated basis matrices (basis functions sampled at
quadrature nodes), and---the one with real arithmetic---\texttt{geom\_data}.
\item[Run-time (apply) stratum.] Data produced afresh on every action
$\vv\mapsto\mat{A}\vv$: the gathered \lstinline[style=l4style]|[E, L]| block, the
evaluated \lstinline[style=l4style]|[E, P, C]| values, the weighted values, the
integrated contributions. These are transient---born inside one application and
discarded when it returns.
\end{description}

The geometry factor is the star of the build-time stratum. Computing it requires
real work: at each quadrature point of each element, form the Jacobian $J$ of the
map from the reference element to the physical element (\cref{ch:refelement}),
combine it into the metric the weak-form term needs---$|J|$ for a mass term,
$J^{-\mathsf T}J^{-1}|J|$ for a gradient term---multiply by the quadrature weight
$w$, and fold in the material coefficient $Q$. The result is one small block per
point, stored in \texttt{geom\_data}. \Cref{fig:mf-strata} draws the split: a
setup pass runs the build once and produces \texttt{geom\_data}; thereafter every
operator application, of which a Krylov solve performs hundreds, simply reads it.

\begin{figure}[t]
\centering
\begin{tikzpicture}[node distance=7mm and 12mm]
  % --- build-time row (gold) ---
  \node[data, fill=simBgGold, draw=simGold, text width=20mm] (mesh)
    {mesh nodes\\$+$ quad rule};
  \node[stage, fill=simBgGold, draw=simGold, right=of mesh, text width=26mm]
    (build) {\textbf{setup pass}\\\footnotesize $J,\;|J|,\;Q,\;w$\\(\texttt{geom\_factor\_build})};
  \node[data, fill=simBgGold, draw=simGold, right=of build, text width=22mm]
    (gd) {\texttt{geom\_data}\\\texttt{[E, P, G]}};
  \draw[flow, simGold] (mesh) -- (build);
  \draw[flow, simGold] (build) -- (gd);
  \node[font=\scriptsize\bfseries, simGold, left=1mm of mesh, rotate=90, anchor=south]
    {build-time};
  \node[font=\scriptsize\itshape, simGray, below=1mm of build]
    {runs \emph{once} per (mesh, order)};

  % --- run-time row (blue): many applies reading geom_data ---
  \node[stage, below=18mm of build, text width=30mm] (apply1)
    {apply $\mat{A}\vv_1$\\\footnotesize $G^{\mathsf T}B^{\mathsf T}\,D\,B\,G$};
  \node[stage, right=5mm of apply1, text width=14mm] (apply2)
    {apply\\$\mat{A}\vv_2$};
  \node[font=\large, right=2mm of apply2] (dots) {$\cdots$};
  \node[stage, right=2mm of dots, text width=14mm] (applyk)
    {apply\\$\mat{A}\vv_k$};
  \node[font=\scriptsize\bfseries, simBlue, left=1mm of apply1, rotate=90, anchor=south]
    {run-time};
  \node[font=\scriptsize\itshape, simGray, below=1mm of apply1, xshift=12mm]
    {runs \emph{every} Krylov iteration --- reads \texttt{geom\_data}, never rebuilds it};
  % geom_data feeds every apply
  \draw[flow, simGold] (gd.south) -- ++(0,-0.4) -| (apply1.north);
  \draw[flow, simGold] (gd.south) -- ++(0,-0.4) -| (apply2.north);
  \draw[flow, simGold] (gd.south) -- ++(0,-0.4) -| (applyk.north);
\end{tikzpicture}
\caption[Build-time versus run-time stratification]{The geometry factor is
computed \emph{once}. A setup pass (gold, top) consumes the mesh nodes and the
quadrature rule and produces \texttt{geom\_data}, a
\lstinline[style=l4style]|[E, P, G]| tensor holding the Jacobian metric,
material coefficient, and quadrature weight at every point. Every subsequent
operator application (blue, bottom)---and an iterative solve performs
hundreds---simply \emph{reads} \texttt{geom\_data} in its weight stage $D$. The
expensive geometry work is amortized across the whole solve. Only when the mesh
changes (e.g.\ adaptive refinement, \cref{ch:amr}) is the setup pass rerun.}
\label{fig:mf-strata}
\end{figure}

\begin{keyidea}
The matrix-free operator carries two strata. The \emph{build-time} stratum---the
restriction map, the tabulated basis, and the geometry factor
\texttt{geom\_data}---is computed once per $(\text{mesh},\text{order})$ and
reused. The \emph{run-time} stratum---the gathered, evaluated, weighted,
integrated tensors---is rebuilt on every application. Folding the coefficient,
Jacobian metric, and quadrature weight into one precomputed \texttt{geom\_data}
factor is what makes the run-time weight stage $D$ a single pointwise multiply.
The same stratification organizes the whole simulator (\cref{ch:strata}).
\end{keyidea}

\subsection{The diagonal $D$ is embarrassingly parallel}
\label{sec:mf-diagonal}

With \texttt{geom\_data} precomputed, the run-time weight stage $D$ is the
simplest object in the entire finite-element machine: a pointwise multiply.
\Cref{fig:mf-diagonal} makes the structure explicit. Lay the
\lstinline[style=l4style]|[E, P, C]| evaluated field and the
\lstinline[style=l4style]|[E, P, G]| geometry factor side by side; at each point
$(e,p)$, multiply the field block by the geometry block, independently. No point
talks to any other point. There is no reduction, no recurrence, no
order-dependence: the output is a function of the input \emph{at the same index}
and nothing else.

\begin{figure}[t]
\centering
\begin{tikzpicture}[scale=1.0]
  % grid of quad points (E x P flattened), each an independent multiply
  \def\nx{5}\def\ny{3}
  \foreach \i in {0,...,4}{
    \foreach \j in {0,...,2}{
      \node[draw=simGreen, fill=simBgGreen, rounded corners=1pt,
            minimum size=6.5mm, font=\scriptsize] (c\i\j) at (\i*1.0, \j*0.85) {$u_{\i\j}$};
    }
  }
  % geom blocks to the right, one per point (drawn as a thin gold strip overlay)
  \foreach \i in {0,...,4}{
    \foreach \j in {0,...,2}{
      \node[font=\scriptsize, simGold] at (\i*1.0+0.34, \j*0.85+0.34) {$\odot g$};
    }
  }
  % axis labels
  \node[font=\scriptsize\itshape, simRust] at (-0.85, 0.85) {$(e,p)$ index};
  \node[font=\scriptsize, simInk] at (2.0,-0.85) {$P$ points $\times$ $E$ elements (flattened)};
  % big brace: all independent
  \draw[decorate, decoration={brace, amplitude=5pt}, simBlue]
    (-0.45, 2.05) -- (4.45, 2.05)
    node[midway, above=5pt, font=\footnotesize, simBlue]
    {every cell independent --- no coupling, any order, fully parallel};
\end{tikzpicture}
\caption[The quad-point weight as a parallel diagonal]{The weight stage
$D$ visualized over the flattened $(e,p)$ index. Each cell is one quadrature
point's evaluated value $u_{ep}$, multiplied ($\odot$) by its own precomputed
geometry block $g_{ep}$ from \texttt{geom\_data}. No cell depends on any other:
$D$ is \emph{diagonal} in the $(E,P)$ index, the embarrassingly parallel core of
the operator. This is the per-point ``lift'' the implementation view consumes
directly---one tensor, one elementwise multiply, mapped across the whole point
set at once.}
\label{fig:mf-diagonal}
\end{figure}

This is what we mean when we call $D$ the \emph{diagonal} of the operator. The
basis stages $B_{\diffop}$ and $B_{\diffop}^{\mathsf T}$ couple the $L$ local dofs
to the $P$ quadrature points (a dense small contraction); the gather and scatter
couple elements through shared dofs. But sandwiched between them, $D$ couples
nothing---it is a diagonal scaling in the quadrature-point basis. A diagonal map
is the friendliest possible object for a parallel machine, and it is where the
material physics ($Q$) and the geometry ($J$) actually enter the operator. Every
bit of what makes this term a \emph{mass} term rather than a \emph{stiffness}
term, this material rather than that one, this curved element rather than a
straight one, is packed into the diagonal $D$.

\section{Worked example: a 1-D mass operator, two elements}
\label{sec:mf-wex1}

Nothing convinces like tracing a real vector through the pipeline and recovering
the answer the assembled matrix would have given. We take the smallest nontrivial
case: the mass operator $\mat{M}$ ($\diffop=\mathrm{Id}$, weight $Q=1$) on a 1-D
mesh of two linear elements, evaluated by a two-point quadrature exact for the
quadratics that arise.

\begin{workedexample}{Matrix-free $\mat{M}\vv$ on two 1-D elements}{mf-mass}
\textbf{Setup.} The domain $[0,2]$ is split into two elements,
$e_0=[0,1]$ and $e_1=[1,2]$, each of width $h=1$. We use linear ($p=1$)
nodal basis functions, so there are $N=3$ global dofs at nodes
$x_0=0,\,x_1=1,\,x_2=2$, and $L=2$ local dofs per element. The node $x_1$ is
\emph{shared}: it is local dof $1$ of $e_0$ and local dof $0$ of $e_1$. Take
$P=2$ Gauss points per element. We compute the action $\mat{M}\vv$ on
\[
  \vv \;=\; (v_0,v_1,v_2) \;=\; (1,\,2,\,3) .
\]

\textbf{Build-time pieces.} On the reference element $[0,1]$ the linear basis is
$\phi_0(\xi)=1-\xi,\ \phi_1(\xi)=\xi$. The two Gauss points are
$\xi_{0,1}=\tfrac12\mp\tfrac{1}{2\sqrt3}$ with weights $w=\tfrac12$ each (scaled
for $[0,1]$). The basis-evaluation matrix $B$ (rows = quad points, cols = local
dofs) and the geometry factor are
\[
  B=\begin{pmatrix}\phi_0(\xi_0)&\phi_1(\xi_0)\\[2pt]\phi_0(\xi_1)&\phi_1(\xi_1)\end{pmatrix}
   =\begin{pmatrix}0.7887&0.2113\\[2pt]0.2113&0.7887\end{pmatrix},
  \qquad
  g_{ep}=\underbrace{|J|}_{=\,h=1}\cdot \underbrace{w}_{=\,1/2}\cdot \underbrace{Q}_{=1}
        =\tfrac12 .
\]
Both elements are identical (uniform mesh), so $B$ and $g=\tfrac12$ are the same
for $e_0$ and $e_1$. These are computed \emph{once}; the four stages below run on
$\vv$.

\tcblower
\textbf{Stage 1 --- gather $G$.} Read each element's two dofs from $\vv=(1,2,3)$:
\[
  e_0:\ (v_0,v_1)=(1,2),\qquad e_1:\ (v_1,v_2)=(2,3)
  \quad\Rightarrow\quad
  \texttt{[E, L]}=\begin{pmatrix}1&2\\2&3\end{pmatrix}.
\]
Note $v_1=2$ is copied into \emph{both} rows---the shared dof.

\textbf{Stage 2 --- evaluate $B$.} Contract each element's local dofs against
$B$ to get values at the two quad points (here $C=1$, scalar values):
\[
  e_0:\ B\begin{pmatrix}1\\2\end{pmatrix}=\begin{pmatrix}1.211\\1.789\end{pmatrix},
  \qquad
  e_1:\ B\begin{pmatrix}2\\3\end{pmatrix}=\begin{pmatrix}2.211\\2.789\end{pmatrix}.
\]
Shape is now \lstinline[style=l4style]|[E, P, C]| $=[2,2,1]$.

\textbf{Stage 3 --- weight $D$.} Multiply every quad-point value by $g=\tfrac12$:
\[
  e_0:\ (0.6057,\,0.8943),\qquad e_1:\ (1.1057,\,1.3943).
\]

\textbf{Stage 4 --- integrate $B^{\mathsf T}$.} Apply $B^{\mathsf T}$ (the
transpose) to return to local dofs:
\[
  e_0:\ B^{\mathsf T}\!\begin{pmatrix}0.6057\\0.8943\end{pmatrix}
       =\begin{pmatrix}0.6667\\0.8333\end{pmatrix},
  \quad
  e_1:\ B^{\mathsf T}\!\begin{pmatrix}1.1057\\1.3943\end{pmatrix}
       =\begin{pmatrix}1.1667\\1.3333\end{pmatrix}.
\]

\textbf{Stage 5 --- scatter-add $G^{\mathsf T}$.} Sum the element-local
contributions back into the $3$ global slots, \emph{adding} at the shared node
$x_1$:
\[
  \mat{M}\vv=\Big(\,0.6667,\ \underbrace{0.8333+1.1667}_{\text{shared }x_1},\ 1.3333\,\Big)
            =(0.6667,\ 2.0000,\ 1.3333).
\]

\textbf{Check against the assembled matrix.} The consistent linear mass matrix on
this uniform two-element mesh is the classic
\[
  \begin{aligned}
  \mat{M}&=\frac{h}{6}\begin{pmatrix}2&1&0\\1&4&1\\0&1&2\end{pmatrix}
        =\frac16\begin{pmatrix}2&1&0\\1&4&1\\0&1&2\end{pmatrix}, \\
  \mat{M}\vv&=\frac16\begin{pmatrix}2{\cdot}1+1{\cdot}2\\1{\cdot}1+4{\cdot}2+1{\cdot}3\\1{\cdot}2+2{\cdot}3\end{pmatrix}
            =\frac16\begin{pmatrix}4\\12\\8\end{pmatrix}
            =(0.6667,\,2.0000,\,1.3333).
  \end{aligned}
\]
The two agree exactly. The matrix-free pipeline reproduced $\mat{M}\vv$ to the
last digit---without ever forming the $3\times3$ matrix.
\end{workedexample}

\begin{remark}
Trace what happened at the shared node $x_1$. The gather $G$ \emph{copied} $v_1$
into both elements (stage 1); the scatter $G^{\mathsf T}$ \emph{summed} the two
elements' contributions back (stage 5). That copy-out/sum-back pair is the entire
content of finite-element assembly, and here it is, applied to a vector instead
of stamped into a matrix. The middle three stages never knew the node was
shared---they ran element by element, blissfully parallel.
\end{remark}

\section{Sum-factorization: the same answer, far cheaper}
\label{sec:mf-sumfac}

The worked example was 1-D, where the basis evaluation $B$ is already a small
dense matrix and there is nothing to optimize. In higher dimensions on
tensor-product elements, the basis stage hides a beautiful and important
efficiency, \emph{sum-factorization}, which is what makes high-order
matrix-free competitive. It is a \emph{transparent} performance trick in the
sense of \cref{ch:bigidea}: it computes exactly the same result, just by
reordering the contraction.

Consider a 2-D tensor-product ($Q_p$) element of order $p$, with
$L=(p+1)^2$ local dofs arranged on a $(p+1)\times(p+1)$ grid, and $P=(p+1)^2$
quadrature points likewise on a grid. The naive basis evaluation treats $B$ as a
dense $P\times L$ matrix: contracting one element's dofs against it costs
\[
  \text{naive cost} \;=\; P\cdot L \;=\; (p+1)^2\cdot(p+1)^2 \;=\; (p+1)^4
\]
multiply-adds per element. But the 2-D basis is a \emph{product} of
one-dimensional bases: $\phi_{ij}(\xi,\eta)=\phi_i(\xi)\,\phi_j(\eta)$. That
product structure lets us evaluate one dimension at a time. First contract the
$\eta$-direction dofs against the 1-D basis (a $(p+1)\times(p+1)$ matrix applied
along one axis), then contract the $\xi$-direction. Each of the two
sweeps costs $(p+1)$ (the 1-D contraction) times $(p+1)^2$ (the other axis carried
along), so
\[
  \text{sum-factored cost} \;=\; d\cdot(p+1)\cdot(p+1)^2 \;=\; d\,(p+1)^{d+1}
  \;=\; 2\,(p+1)^3 \quad (d=2).
\]
The general $d$-dimensional statement is the move from $O(P\cdot L)=O\!\big((p+1)^{2d}\big)$
to $O\!\big(d\,(p+1)^{d+1}\big)$---in the notation of the operator counts, a drop
from $O(P\cdot L)$ to $O\!\big(d\cdot P^{1/d}\cdot L\big)$, since a single axis of
the quad-point grid has length $P^{1/d}=(p+1)$.

\Cref{fig:mf-sumfac} draws the reordering for $d=2$, and \cref{tab:mf-sumfac}
tabulates the savings. The win grows with order: at $p=8$ in 3-D it is the
difference between $(p+1)^6\approx\num{531441}$ and
$3\cdot(p+1)^4\approx\num{19683}$ operations per element---a factor of nearly
$30$, and it widens as $p$ climbs.

\begin{figure}[t]
\centering
\begin{tikzpicture}[node distance=6mm and 14mm]
  % naive: one big dense contraction L -> P
  \node[data, text width=15mm] (din) {dofs\\$(p{+}1)^2$};
  \node[opbox, fill=simBgRust, draw=simRust, right=of din, text width=24mm]
    (naive) {dense $P\times L$\\\footnotesize $(p{+}1)^4$ ops};
  \node[data, right=of naive, text width=16mm] (dpts) {pts\\$(p{+}1)^2$};
  \draw[flow] (din) -- (naive);
  \draw[flow] (naive) -- (dpts);
  \node[font=\scriptsize\bfseries, simRust, left=1mm of din, rotate=90, anchor=south]
    {naive};

  % factored: two 1-D sweeps
  \node[data, below=14mm of din, text width=15mm] (fin) {dofs\\$(p{+}1)^2$};
  \node[opbox, right=of fin, text width=20mm] (s1)
    {1-D sweep $\eta$\\\footnotesize $(p{+}1)^3$ ops};
  \node[opbox, right=of s1, text width=20mm] (s2)
    {1-D sweep $\xi$\\\footnotesize $(p{+}1)^3$ ops};
  \node[data, right=of s2, text width=16mm] (fpts) {pts\\$(p{+}1)^2$};
  \draw[flow] (fin) -- (s1);
  \draw[flow] (s1) -- node[above, font=\scriptsize, simGray] {partial} (s2);
  \draw[flow] (s2) -- (fpts);
  \node[font=\scriptsize\bfseries, simBlue, left=1mm of fin, rotate=90, anchor=south]
    {factored};
\end{tikzpicture}
\caption[Sum-factorization on a 2-D element]{Basis evaluation on a 2-D
tensor-product element, two ways. \emph{Top (rust):} the naive contraction treats
the basis as one dense $P\times L$ matrix, costing $(p+1)^4$ operations per
element. \emph{Bottom (blue):} sum-factorization applies the one-dimensional
basis along each axis in turn---a sweep in $\eta$, then a sweep in $\xi$---each
costing $(p+1)^3$, for $2(p+1)^3$ total. The two produce \emph{identical}
results; only the contraction order differs. The saving is the move from
$O(P\cdot L)$ to $O(d\,P^{1/d}L)$, and it grows with order.}
\label{fig:mf-sumfac}
\end{figure}

\begin{table}[t]
\centering
\caption[Sum-factorization operation counts]{Multiply-adds per element for the
basis-evaluation stage, naive versus sum-factored, on a tensor-product element of
order $p$. ``Naive'' is the dense $P\times L$ contraction, $(p+1)^{2d}$;
``factored'' is $d\,(p+1)^{d+1}$. The speedup is unbounded in $p$. (Values for
$d=2$; the 3-D column at $p=8$ illustrates how steep the win becomes.)}
\label{tab:mf-sumfac}
\small
\setlength{\tabcolsep}{8pt}
\renewcommand{\arraystretch}{1.2}
\begin{tabular}{@{}cccc@{}}
\toprule
order $p$ & naive $(p{+}1)^4$ & factored $2(p{+}1)^3$ & speedup \\
\midrule
1 & 16     & 16    & $1.0\times$ \\
2 & 81     & 54    & $1.5\times$ \\
4 & 625    & 250   & $2.5\times$ \\
8 & 6561   & 1458  & $4.5\times$ \\
\midrule
\multicolumn{4}{@{}l@{}}{\footnotesize 3-D at $p=8$: naive $9^6=\num{531441}$
vs.\ factored $3\cdot 9^4=\num{19683}$ \quad($27\times$)} \\
\bottomrule
\end{tabular}
\end{table}

\begin{workedexample}{Counting operations on a $Q_2$ element}{mf-sumfac}
\textbf{Setup.} Take a 2-D second-order ($Q_2$, $p=2$) tensor-product element.
The local dofs sit on a $3\times3$ node grid, so $L=(p+1)^2=9$; we use a $3\times3$
Gauss grid, $P=9$. We count the multiply-adds to evaluate the field values at all
quad points (the basis stage $B$, mode = interp).

\tcblower
\textbf{Naive (dense $P\times L$).} One element's $L=9$ dofs contract against the
dense $9\times9$ basis matrix:
\[
  P\cdot L \;=\; 9\cdot 9 \;=\; 81 \ \text{multiply-adds}.
\]

\textbf{Factored (two 1-D sweeps).} The $Q_2$ basis is
$\phi_{ij}(\xi,\eta)=\phi_i(\xi)\phi_j(\eta)$ with $\phi_i,\phi_j$ the three 1-D
quadratic Lagrange functions.
\emph{Sweep 1 (along $\eta$):} for each of the $3$ values of $\xi$-index, contract
the $3$ dofs in $\eta$ against the $3\times3$ 1-D basis, producing partial values
at the $3$ $\eta$-quad-points: $3\ (\xi\text{ cols})\times 3\times 3 = 27$.
\emph{Sweep 2 (along $\xi$):} for each of the $3$ $\eta$-quad-points, contract the
$3$ partials in $\xi$ against the $3\times3$ 1-D basis: $3\times 3\times3 = 27$.
\[
  \text{total} \;=\; 27 + 27 \;=\; 54 \ \text{multiply-adds}.
\]

\textbf{Comparison.} $81$ versus $54$: a $1.5\times$ saving at $p=2$, matching the
$d(p+1)^{d+1}=2\cdot 27=54$ formula. The two compute \emph{bit-for-bit the same
quad-point values}---sum-factorization only reorders the sums. The modest factor
here grows fast: at $p=4$ it is $625$ vs.\ $250$ ($2.5\times$), and in 3-D the
exponent gap ($2d$ vs.\ $d+1$) makes it decisive.
\end{workedexample}

\begin{pitfall}
Sum-factorization is \emph{transparent}: same result, faster. Do not confuse it
with a \emph{load-bearing} numerical trick (\cref{ch:bigidea})---it does not
change the answer, the conditioning, or the rounding behavior in any way the
algorithm depends on. In the spec we therefore record it as a one-line note on
the basis stage, not as a separate algebraic claim. The reason to know about it
anyway is that it is \emph{why} high-order matrix-free is fast enough to be the
default: without it, the basis stage would reintroduce the very $O(p^{2d})$ cost
the matrix-free form was trying to escape.
\end{pitfall}

\section{The implementation view: contractions all the way down}
\label{sec:mf-impl-view}

Step back from the details and look at what \cref{eq:mf-chain} \emph{is}. The
finite-element operator $\mat{A}$---the discretization of Maxwell's weak form, the
object every solver in the book applies---has turned out to be a composition of
tensor contractions over named axes. Gather is an indexed copy. Evaluate and
integrate are dense contractions of an $[E,L]$ tensor against a tabulated basis,
yielding $[E,P,C]$. Weight is an elementwise multiply of $[E,P,C]$ against
$[E,P,G]$. Scatter is an indexed sum. There is no sparse-matrix data structure
anywhere in sight---no row pointers, no column indices, no stored entries. There
are only tensors and contractions.

This is the destination Part~I promised. \Cref{ch:bigidea} argued that a field
\emph{is} a tensor and that a simulation is a chain of pure transformations of
immutable tensors---and that this description is exactly the form a GPU or tensor
accelerator wants: whole-tensor operations over named axes, no aliasing, no
ordering hazards. The matrix-free operator is that thesis made fully concrete.
Every stage of \cref{fig:mf-pipeline} is a pure function from tensors to tensors;
the whole operator is their composition; and the shapes
(\lstinline[style=l4style]|[E, L]|, \lstinline[style=l4style]|[E, P, C]|,
\lstinline[style=l4style]|[E, P, G]|) are precisely what a backend schedules onto
its parallel units. The assembled-matrix form, by contrast, is a CPU-era data
structure---irregular, pointer-chasing, bandwidth-bound---that a tensor backend
would have to be coaxed into. Matrix-free is the form the accelerator was built
for.

\begin{keyidea}
The matrix-free operator \emph{is} the implementation view of \cref{ch:bigidea},
realized for finite elements. A field is a tensor; the operator is a composition
of contractions over named axes; the whole thing maps onto a tensor backend with
no sparse data structure at all. When the book's final part renders the
synthesized library, the matrix-free operator is what the
``\,\lstinline[style=l4style]|apply A v|\,'' of every solver expands into.
\end{keyidea}

\subsection{The innermost kernel is a library boundary}
\label{sec:mf-extern}

There is one honest qualification. We have described the five stages as pure
tensor contractions, and at the mathematical level they are. But the innermost
pointwise kernel---the per-quadrature-point weight stage $D$, fused with the basis
contractions for efficiency---is, in the real Palace implementation, handed to a
specialized accelerator library (libCEED) that owns the device-level details:
which loops to fuse, how to tile the quadrature points, how to map them onto warps
or wavefronts. The book does not re-derive that kernel; it treats it as an
\emph{opaque library boundary}. In the calculus, such a boundary is written with
the marker \lstinline[style=l4style]|#extern|: a node whose \emph{contract} (its
shape signature and semantics) we state precisely, but whose \emph{body} we
delegate to the library.

\begin{lstlisting}[style=l4style]
-- the innermost element-quadrature kernel, as a library boundary:
-- the contract is ours; the device implementation is libCEED's.
quad_kernel :: GeomData -> Tensor[E, P, C] -> Tensor[E, P, C]
#extern quad_kernel    -- shape + semantics specified; body delegated
\end{lstlisting}

\noindent This is the same \lstinline[style=l4style]|#extern| node that appears,
violet and dashed, at the center of \cref{fig:mf-pipeline}. Two facts make the
boundary clean. First, its contract is exactly the pointwise diagonal of
\cref{sec:mf-diagonal}: a shape-preserving elementwise multiply against
\texttt{geom\_data}, with semantics we have fully specified. Second, because the
contract is so simple and so completely characterized, we can reason about the
\emph{whole} operator---its symmetry, its action, the worked examples
above---without opening the kernel at all. The library owns the machine; we own
the mathematics, and the two meet at a narrow, well-typed interface.

\begin{notationbox}
The \lstinline[style=l4style]|#extern| marker recurs wherever the spec touches an
opaque accelerator or solver library: the libCEED quadrature kernel here, the
SLEPc eigensolve in \cref{ch:eigenvalues}, the sparse triangular solves inside a
smoother in \cref{ch:smoothers}. In each case the pattern is the same---a
precisely specified \emph{contract} (signature and semantics) standing in for a
\emph{body} we delegate. The discipline is to make the boundary so well-typed that
nothing above it needs to look inside.
\end{notationbox}

\section{When matrix-free is not the answer}
\label{sec:mf-tradeoffs}

Matrix-free is the right default at moderate-to-high order on accelerators, but it
is a trade, not a free lunch, and it pays to name the other side of the ledger.
\Cref{fig:mf-vs-assembled} sets the two representations side by side: the same
linear map, stored two ways. The assembled form holds an explicit sparse data
structure---rows, columns, stored entries---and applies in one multiply; the
matrix-free form holds only the build-time tables and recomputes the action
through the five-stage pipeline each time. The choice between them is the subject
of this section.

\begin{figure}[t]
\centering
\begin{tikzpicture}[node distance=5mm and 10mm]
  % ---- LEFT: assembled (stored sparse matrix) ----
  \begin{scope}
    \node[font=\bfseries\small, simRust] at (0,2.1) {Assembled};
    % a sparse matrix glyph
    \draw[simRust, thick, fill=simBgRust, rounded corners=1pt] (-1.0,-1.1) rectangle (1.0,1.1);
    \foreach \i/\j in {0/0,1/0,0/1,1/1,2/2,3/2,2/3,3/3,1/2,2/1,4/4,3/4,4/3,0/4,4/0}{
      \fill[simRust] ({-0.9+0.38*\j},{0.9-0.38*\i}) rectangle ({-0.66+0.38*\j},{0.66-0.38*\i});}
    \node[font=\scriptsize\itshape, simGray, below=1mm of {(0,-1.1)}] (la) at (0,-1.45)
      {stored entries $A_{ij}$};
    \node[data, fill=simBgRust, draw=simRust, below=2mm of la, text width=30mm]
      (lapp) {apply: one sparse multiply\\\footnotesize cheap per-apply, $O(p^{2d})$ memory};
  \end{scope}
  % ---- RIGHT: matrix-free (pipeline, nothing stored) ----
  \begin{scope}[shift={(6.4,0)}]
    \node[font=\bfseries\small, simBlue] at (0,2.1) {Matrix-free};
    % mini pipeline glyph
    \foreach \k/\lab in {0/$G$,1/$B$,2/$D$,3/$B^{\mathsf T}$,4/$G^{\mathsf T}$}{
      \node[draw=simBlue, fill=simBgBlue, rounded corners=1pt, minimum size=4.5mm,
            font=\scriptsize] (p\k) at ({-1.4+0.7*\k},0.6) {\lab};}
    \foreach \k in {0,1,2,3}{\pgfmathtruncatemacro{\next}{\k+1}\draw[flow] (p\k)--(p\next);}
    \node[font=\scriptsize\itshape, simGray] at (0,-0.5) {nothing stored but build-time tables};
    \node[data, fill=simBgBlue, draw=simBlue, text width=30mm] (rapp) at (0,-1.55)
      {apply: rerun the pipeline\\\footnotesize more arithmetic, $O(p)$ memory};
  \end{scope}
  % ---- center: same operator ----
  \node[font=\scriptsize\bfseries, simGreen] at (3.2,1.4) {same};
  \node[font=\scriptsize\bfseries, simGreen] at (3.2,1.05) {$\mat{A}$};
  \draw[{Stealth[length=2mm]}-{Stealth[length=2mm]}, simGreen, thick] (1.3,0.2) -- (4.5,0.2);
\end{tikzpicture}
\caption[Matrix-free versus assembled, side by side]{Two representations of the
\emph{same} operator $\mat{A}$. \emph{Left (rust):} the assembled form stores an
explicit sparse matrix of entries $A_{ij}$ and applies in a single sparse
multiply---cheap per application but $O(p^{2d})$ memory per element.
\emph{Right (blue):} the matrix-free form stores nothing but the build-time
tables and recomputes the action through the five-stage pipeline on each apply---more
arithmetic but only $O(p)$ memory. They compute the identical map (green); the
choice is an implementation decision driven by order, hardware, and downstream use.}
\label{fig:mf-vs-assembled}
\end{figure}

First, \emph{cost per application is higher}. The assembled matrix, once built,
applies in a single sparse multiply; the matrix-free operator reruns all five
stages every time. When the operator is applied only a handful of times, or when
order is so low that the stored matrix is tiny, assembly wins. The crossover in
\cref{fig:mf-cost} is a memory crossover; there is a separate, lower, arithmetic
crossover.

Second, \emph{preconditioning is harder}. Many classical preconditioners
(incomplete factorizations, algebraic multigrid setup) want to inspect the
operator's entries---exactly what the matrix-free form refuses to provide. This is
not fatal: the geometric multigrid of \cref{ch:multigrid} and the smoothers of
\cref{ch:smoothers} are designed to need only the action, and where an entry-level
view is unavoidable, Palace can materialize a coarse or diagonal approximation on
demand (apply the pipeline to the identity columns and read off the entries). But
it is real friction, and it is why the matrix-free and assembled forms coexist in
the same code rather than one simply replacing the other.

\begin{remark}
Both forms compute the \emph{same operator}---they are two evaluations of the same
linear map, and the worked example confirmed they agree to the last digit. The
choice between them is an implementation decision driven by order, hardware, and
how the operator will be used downstream, not a change to the mathematics. This is
the value-threading lesson of \cref{ch:bigidea} once more: the description is free
to choose the representation the machine likes best, because the answer is
invariant to the choice.
\end{remark}

\summarysection
A Krylov or eigensolver asks an operator for one thing only: its action,
$\vv\mapsto\mat{A}\vv$. The \emph{matrix-free} operator supplies that action
without ever forming or storing $\mat{A}$, which matters because the stored
operator costs $O(p^{2d})$ memory per element and becomes unaffordable at high
order on accelerators. The action factors into five element-local stages,
$\mat{A}\vv = G^{\mathsf T}B_{\diffop}^{\mathsf T}\,D\,B_{\diffop}\,G\,\vv$:
\emph{gather} the global vector to per-element blocks
(\lstinline[style=l4style]|[N]|$\to$\lstinline[style=l4style]|[E, L]|),
\emph{evaluate} the basis at quadrature points
(\lstinline[style=l4style]|[E, L]|$\to$\lstinline[style=l4style]|[E, P, C]|),
\emph{weight} each point by a precomputed geometry/material factor (pointwise on
\lstinline[style=l4style]|[E, P, C]|, the embarrassingly parallel diagonal),
\emph{integrate} back to local dofs, and \emph{scatter-add} to a fresh global
vector. The middle three stages are block-diagonal in the element axis $E$---fully
parallel; all inter-element coupling lives in the gather/scatter ends. The
geometry factor \texttt{geom\_data} is a \emph{build-time} quantity, computed once
per $(\text{mesh},\text{order})$ and reused across the hundreds of applications a
solve performs. On tensor-product elements, \emph{sum-factorization} evaluates the
basis dimension by dimension, dropping the cost from $O(P\cdot L)$ to
$O(d\,P^{1/d}L)$---a transparent reordering, same answer, decisive at high order.
A worked example traced $\mat{M}\vv$ through all five stages on two 1-D elements
and recovered the assembled matrix's answer exactly. Because every stage is a
tensor contraction over named axes with no sparse-matrix structure, the
matrix-free operator \emph{is} the implementation view of \cref{ch:bigidea}---the
form a GPU consumes---with its innermost kernel sealed behind a precisely typed
\lstinline[style=l4style]|#extern| library boundary.

\furtherreading
The matrix-free, sum-factorized, tensor-contraction view of high-order
finite-element operators is the design philosophy of the CEED project; its
co-design survey~\citep{ceed2021} is the entry point to the libraries (libCEED,
MFEM) that realize the five-stage decomposition described here, and to the
performance literature behind sum-factorization. For the underlying
finite-element method---reference elements, tabulated bases, quadrature, and the
element matrices the operator factors---Jin's text on the finite-element method in
electromagnetics~\citep{jin2014} gives the field-theory grounding in the
context of Maxwell's equations specifically. The classical operation-count
argument for sum-factorization, and the spectral-element setting in which it was
first decisive, are developed in any spectral/high-order finite-element reference;
the build-time/run-time stratification it relies on is treated in its own right in
\cref{ch:strata}.

\exercisessection
\begin{enumerate}[nosep]
  \item \textbf{Memory at a glance.} For a 3-D hexahedral element of order $p=4$,
  how many entries does the dense element matrix $\mat{A}^e$ hold? How many
  numbers does the matrix-free form store for that element if it keeps a tabulated
  1-D basis of size $(p+1)^2$ and a geometry block of $G=2+d^2$ components at each
  of $(p+1)^3$ quadrature points? Which is larger, and by what factor?

  \item \textbf{Shapes along the pipeline.} A driven-analysis curl--curl operator
  ($\diffop=\Curl$) runs on a mesh with $E=\num{1000}$ hexahedral elements of
  order $p=2$, with $P=27$ quadrature points per element and $C=3$ curl components.
  Write the tensor shape on each of the five wires of \cref{fig:mf-pipeline},
  from the input \lstinline[style=l4style]|[N]| to the output
  \lstinline[style=l4style]|[N]|, filling in concrete extents for $E$, $L$, $P$,
  $C$. (You will need $L=(p+1)^3$.)

  \item \textbf{The shared-dof bookkeeping.} In
  \cref{wex:mf-mass}, the node $x_1$ received contributions
  $0.8333$ (from $e_0$) and $1.1667$ (from $e_1$), summing to $2.0000$. Explain,
  in terms of $G$ and $G^{\mathsf T}$, why this \emph{sum} is the correct global
  value and not, say, an average. What property of $G^{\mathsf T}$ as the
  transpose of $G$ guarantees it?

  \item \textbf{Counting sum-factorization.} Repeat the operation count of
  \cref{wex:mf-sumfac} for a 3-D $Q_2$ element ($d=3$, $p=2$, so $L=P=27$).
  Give the naive count $P\cdot L$ and the sum-factored count $d\,(p+1)^{d+1}$, and
  report the speedup. Then redo it for $p=4$ and comment on how the speedup scales.

  \item \textbf{The diagonal carries the physics.} Two terms---a mass term
  ($\diffop=\mathrm{Id}$, $Q=\varepsilon$) and a stiffness term ($\diffop=\Grad$,
  $Q=\varepsilon$)---are applied on the \emph{same} mesh. Which of the five stages
  differ between them, and which are identical? In particular, what changes in the
  contents (not the shape) of \texttt{geom\_data}, and what changes in the basis
  mode? Relate your answer to why a single five-stage skeleton serves all six
  analyses of \cref{ch:targets}.

  \item \textbf{Build-time amortization.} A conjugate-gradient solve
  (\cref{ch:cg}) converges in $200$ iterations, each requiring one operator
  application. The build pass that produces \texttt{geom\_data} costs roughly as
  much as $5$ operator applications. What fraction of the total operator-related
  work is the one-time build? Now suppose adaptive refinement (\cref{ch:amr})
  rebuilds \texttt{geom\_data} after every $50$ iterations. How does the fraction
  change, and what does this say about when matrix-free's amortization advantage
  is strongest?

  \item \textbf{Why the solver never asks for an entry.} Conjugate gradients
  (\cref{ch:cg}) forms, each iteration, the quantities $\mat{A}\vp$,
  $\inner{\vp}{\mat{A}\vp}$, and vector updates $\vx\leftarrow\vx+\alpha\vp$.
  Identify which of these touch the operator, and confirm that none of them
  requires an individual entry $A_{ij}$. What would have to change about the
  \emph{solver} for the matrix-free form to become unusable?

  \item \textbf{The \texttt{\#extern} boundary.} The innermost kernel is sealed
  behind \lstinline[style=l4style]|#extern quad_kernel| with the contract
  \lstinline[style=l4style]|GeomData -> Tensor[E, P, C] -> Tensor[E, P, C]|.
  State, in one sentence each, (a) the semantic property of this kernel that lets
  us reason about the symmetry of $\mat{A}=G^{\mathsf T}B^{\mathsf T}DBG$ without
  opening the kernel, and (b) one reason the implementation might want to
  \emph{fuse} this kernel with the adjacent basis stages even though, on paper,
  $D$ is a separate stage.
\end{enumerate}
